
\documentclass{article}
\usepackage{colortbl}
\usepackage{makecell}
\usepackage{multirow}
\usepackage{supertabular}

\begin{document}

\newcounter{utterance}

\centering \large Interaction Transcript for game `hot\_air\_balloon', experiment `air\_balloon\_survival\_en\_reasoning off\_hard', episode 5 with Qwen3{-}32B{-}t0.0.
\vspace{24pt}

{ \footnotesize  \setcounter{utterance}{1}
\setlength{\tabcolsep}{0pt}
\begin{supertabular}{c@{$\;$}|p{.15\linewidth}@{}p{.15\linewidth}p{.15\linewidth}p{.15\linewidth}p{.15\linewidth}p{.15\linewidth}}
    \# & $\;$A & \multicolumn{4}{c}{Game Master} & $\;\:$B\\
    \hline

    \theutterance \stepcounter{utterance}  
    & & \multicolumn{4}{p{0.6\linewidth}}{
        \cellcolor[rgb]{0.9,0.9,0.9}{
            \makecell[{{p{\linewidth}}}]{
                \texttt{\tiny{[P1$\langle$GM]}}
                \texttt{You are participating in a collaborative negotiation game.} \\
\\ 
\texttt{Together with another participant, you must agree on a single set of items that will be kept. Each of you has your own view of how much each item matters to you (importance). You do not know how the other participant values the items. Additionally, you are given the effort each item demands.} \\
\texttt{You may only agree on a set if the total effort of the selected items does not exceed a shared limit:} \\
\\ 
\texttt{LIMIT = 2307} \\
\\ 
\texttt{Here are the individual item effort values:} \\
\\ 
\texttt{Item effort = \{"A21": 147, "B93": 698, "A11": 588, "B33": 43, "C08": 356, "B39": 80, "B96": 95, "C85": 743, "A82": 106, "B31": 308, "B20": 325, "A60": 255, "B09": 276, "C81": 543, "B21": 51\}} \\
\\ 
\texttt{Here is your personal view on the importance of each item:} \\
\\ 
\texttt{Item importance values = \{"A21": 138, "B93": 583, "A11": 65, "B33": 262, "C08": 121, "B39": 508, "B96": 461, "C85": 484, "A82": 668, "B31": 389, "B20": 215, "A60": 97, "B09": 500, "C81": 30, "B21": 400\}} \\
\\ 
\texttt{Goal:} \\
\\ 
\texttt{Your goal is to negotiate a shared set of items that benefits you as much as possible (i.e., maximizes total importance to YOU), while staying within the LIMIT. You are not required to make a PROPOSAL in every message {-} you can simply negotiate as well. All tactics are allowed!} \\
\\ 
\texttt{Interaction Protocol:} \\
\\ 
\texttt{You may only use the following structured formats in a message:} \\
\\ 
\texttt{PROPOSAL: \{'A', 'B', 'C', …\}} \\
\texttt{Propose keeping exactly those items.} \\
\texttt{REFUSE: \{'A', 'B', 'C', …\}} \\
\texttt{Explicitly reject opponent's proposal.} \\
\texttt{ARGUMENT: \{'...'\}} \\
\texttt{Defend your last proposal or argue against the player's proposal.} \\
\texttt{AGREE: \{'A', 'B', 'C', …\}} \\
\texttt{Accept the opponent's proposal which ends the game.} \\
\\ 
\\ 
\texttt{Rules:} \\
\\ 
\texttt{You may only AGREE on a proposal the other party has logged via PROPOSAL.} \\
\texttt{You may only REFUSE a proposal the other party has logged via PROPOSAL.} \\
\texttt{Total effort of any PROPOSAL or AGREE set must be ≤ LIMIT.} \\
\texttt{Do NOT reveal your hidden importance scores.} \\
\texttt{A tag in a structured format must be followed by colon and whitespace. The argument must be a python set containing 0 or more strings.} \\
\texttt{So, it must be of the form TAG: \{...\}} \\
\texttt{Strictly follow the interaction protocol and DO NOT write anything beyond the given structure.} \\
\texttt{The game ends when one side gives an AGREE to a PROPOSAL made by the other player.} \\
\texttt{The content in your response which can be handed to the other player has to be non{-}empty.} \\
\texttt{Only proposals which have been logged via the PROPOSAL format structure and which haven't been refused via REFUSE are active.} \\
\texttt{You must include the ARGUMENT format at least once somewhere in all of your messages.} \\
            }
        }
    }
    & \\ \\

    \theutterance \stepcounter{utterance}  
    & \multicolumn{4}{p{0.6\linewidth}}{
        \cellcolor[rgb]{0.9,0.9,0.9}{
            \makecell[{{p{\linewidth}}}]{
                \texttt{\tiny{[P1$\rangle$GM]}}
                \texttt{<think>} \\
\texttt{Okay, let's see. I need to negotiate a set of items that maximizes my own importance while keeping the total effort under 2307. First, I should calculate the total effort of all items to see if we can take them all. Let me add up the effort values.} \\
\\ 
\texttt{Looking at the item efforts: A21 is 147, B93 is 698, A11 is 588, B33 is 43, C08 is 356, B39 is 80, B96 is 95, C85 is 743, A82 is 106, B31 is 308, B20 is 325, A60 is 255, B09 is 276, C81 is 543, B21 is 51. Let me sum these up step by step.} \\
\\ 
\texttt{Starting with 0:} \\
\texttt{+147 = 147} \\
\texttt{+698 = 845} \\
\texttt{+588 = 1433} \\
\texttt{+43 = 1476} \\
\texttt{+356 = 1832} \\
\texttt{+80 = 1912} \\
\texttt{+95 = 2007} \\
\texttt{+743 = 2750. Oh, that's already over 2307. So taking all items isn't possible. So we need to remove some items to get under the limit.} \\
\\ 
\texttt{The total of all items is 2750, which is 443 over the limit. So we need to remove items whose combined effort is at least 443. Now, I need to figure out which items to remove to minimize the loss of my importance. Since my goal is to maximize my own importance, I should prioritize keeping the items with high importance per unit effort.} \\
\\ 
\texttt{Let me calculate the importance per effort ratio for each item. But maybe a better approach is to sort the items by importance descending and see which ones are the most valuable to me, then see if removing the least important ones can get us under the limit.} \\
\\ 
\texttt{My importance values are:} \\
\texttt{A21:138, B93:583, A11:65, B33:262, C08} \\
            }
        }
    }
    & & \\ \\

    \theutterance \stepcounter{utterance}  
    & & \multicolumn{4}{p{0.6\linewidth}}{
        \cellcolor[rgb]{0.9,0.9,0.9}{
            \makecell[{{p{\linewidth}}}]{
                \texttt{\tiny{[P1$\langle$GM]}}
                \texttt{Your response did not contain an argument. You must include the structured format ARGUMENT: \{...\} somewhere in your response. Try again} \\
            }
        }
    }
    & \\ \\

    \theutterance \stepcounter{utterance}  
    & \multicolumn{4}{p{0.6\linewidth}}{
        \cellcolor[rgb]{0.9,0.9,0.9}{
            \makecell[{{p{\linewidth}}}]{
                \texttt{\tiny{[P1$\rangle$GM]}}
                \texttt{<think>} \\
\texttt{Okay, let me try to figure this out. The user wants me to make a proposal that maximizes my own importance while staying under the effort limit. First, I need to calculate the total effort of all items to see if we can take them all. Let me add up the effort values.} \\
\\ 
\texttt{Looking at the item efforts: A21 is 147, B93 is 698, A11 is 588, B33 is 43, C08 is 356, B39 is 80, B96 is 95, C85 is 743, A82 is 106, B31 is 308, B20 is 325, A60 is 255, B09 is 276, C81 is 543, B21 is 51. Let me sum these up step by step.} \\
\\ 
\texttt{Starting with 0:} \\
\texttt{+147 = 147} \\
\texttt{+698 = 845} \\
\texttt{+588 = 1433} \\
\texttt{+43 = 1476} \\
\texttt{+356 = 1832} \\
\texttt{+80 = 1912} \\
\texttt{+95 = 2007} \\
\texttt{+743 = 2750. Oh, that's already over 2307. So taking all items isn't possible. The total is 2750, which is 443 over the limit. So we need to remove items whose combined effort is at least 443.} \\
\\ 
\texttt{Now, I need to decide which items to remove to minimize the loss of my importance. My goal is to maximize my own importance, so I should prioritize keeping the items with the highest importance per unit effort. Let me calculate the importance per effort ratio for each item.} \\
\\ 
\texttt{My importance values are:} \\
\texttt{A21:138, B93:583, A11:65, B33:262, C08:121, B39:508, B96:461, C85:484, A82:668, B31:389, B20:215, A6} \\
            }
        }
    }
    & & \\ \\

    \theutterance \stepcounter{utterance}  
    & & \multicolumn{4}{p{0.6\linewidth}}{
        \cellcolor[rgb]{0.9,0.9,0.9}{
            \makecell[{{p{\linewidth}}}]{
                \texttt{\tiny{[P1$\langle$GM]}}
                \texttt{Your response did not contain an argument. You must include the structured format ARGUMENT: \{...\} somewhere in your response. Try again} \\
            }
        }
    }
    & \\ \\

    \theutterance \stepcounter{utterance}  
    & \multicolumn{4}{p{0.6\linewidth}}{
        \cellcolor[rgb]{0.9,0.9,0.9}{
            \makecell[{{p{\linewidth}}}]{
                \texttt{\tiny{[P1$\rangle$GM]}}
                \texttt{<think>} \\
\texttt{Okay, let me try to figure this out. The user wants me to make a proposal that maximizes my own importance while staying under the effort limit. First, I need to calculate the total effort of all items to see if we can take them all. Let me add up the effort values.} \\
\\ 
\texttt{Looking at the item efforts: A21 is 147, B93 is 698, A11 is 588, B33 is 43, C08 is 356, B39 is 80, B96 is 95, C85 is 743, A82 is 106, B31 is 308, B20 is 325, A60 is 255, B09 is 276, C81 is 543, B21 is 51. Let me sum these up step by step.} \\
\\ 
\texttt{Starting with 0:} \\
\texttt{+147 = 147} \\
\texttt{+698 = 845} \\
\texttt{+588 = 1433} \\
\texttt{+43 = 1476} \\
\texttt{+356 = 1832} \\
\texttt{+80 = 1912} \\
\texttt{+95 = 2007} \\
\texttt{+743 = 2750. Oh, that's already over 2307. So taking all items isn't possible. The total is 2750, which is 443 over the limit. So we need to remove items whose combined effort is at least 443.} \\
\\ 
\texttt{Now, I need to decide which items to remove to minimize the loss of my importance. My goal is to maximize my own importance, so I should prioritize keeping the items with the highest importance per unit effort. Let me calculate the importance per effort ratio for each item.} \\
\\ 
\texttt{My importance values are:} \\
\texttt{A21:138, B93:583, A11:65, B33:262, C08:121, B39:508, B96:461, C85:484, A82:668, B31:389, B20:215, A6} \\
            }
        }
    }
    & & \\ \\

    \theutterance \stepcounter{utterance}  
    & & & \multicolumn{2}{p{0.3\linewidth}}{
        \cellcolor[rgb]{0.9,0.9,0.9}{
            \makecell[{{p{\linewidth}}}]{
                \texttt{\tiny{[GM$|$GM]}}
                \texttt{abort} \\
            }
        }
    }
    & & \\ \\

    \theutterance \stepcounter{utterance}  
    & & & \multicolumn{2}{p{0.3\linewidth}}{
        \cellcolor[rgb]{0.9,0.9,0.9}{
            \makecell[{{p{\linewidth}}}]{
                \texttt{\tiny{[GM$|$GM]}}
                \texttt{end game} \\
            }
        }
    }
    & & \\ \\

\end{supertabular}
}

\end{document}
